\documentclass[11pt]{article}

\usepackage[margin=1in]{geometry}
\usepackage{amsmath, amssymb}
\usepackage{booktabs}

\title{Consumption Tilt Across Periods with Different Equivalence Scales}
\author{Divorce branch -- modeling note}
\date{\today}

\begin{document}
\sloppy
\maketitle

\section{Setup}

Two periods with equivalence scales $s_1 = \sqrt{\text{hh\_size}_1}$,
$s_2 = \sqrt{\text{hh\_size}_2}$ and utility
$u(C_t) = \big((C_t/s_t)^{1-\mu} - 1\big)/(1-\mu)$ (no household-size
multiplier in front, per \texttt{utility\_functions\_add.py}). Maximizing
$u(C_1) + \beta u(C_2)$ subject to $C_1 + C_2/R = W$ gives the Euler
equation
\begin{equation}
  \frac{C_2}{C_1} = (\beta R)^{1/\mu} \left(\frac{s_1}{s_2}\right)^{(1-\mu)/\mu}.
\end{equation}
For comparison, with the $\text{hh\_size}_t = s_t^2$ multiplier back in
front ($\tilde u(C_t) = \text{hh\_size}_t \cdot u(C_t)$, the pre-fix
version), the same steps give
\begin{equation}
  \frac{C_2}{C_1} = (\beta R)^{1/\mu} \left(\frac{s_1}{s_2}\right)^{-(\mu+1)/\mu}.
\end{equation}

\section{Result}

The sign of the exponent on $(s_1/s_2)$ determines which period the tilt
favors. Without the multiplier the sign depends on $\mu$; with it, the
exponent $-(\mu+1)/\mu$ is negative for every $\mu>0$, so it \emph{never}
flips -- it always favors the larger household. Table~\ref{tab:tilt} gives
the direction for each regime, independent of which period happens to be
larger.

\begin{table}[h]
\centering
\begin{tabular}{@{}llll@{}}
\toprule
$\mu$ & Pre-fix (with $\text{hh\_size}_t$ multiplier) & Post-fix (no multiplier) & example ($s_1>s_2$) \\
\midrule
$\mu > 1$ & larger scale & larger scale & period 1 (partnered) \\
$\mu = 1$ & larger scale & no tilt, $C_2/C_1=\beta R$ & period 1 (partnered) \\
$\mu < 1$ & larger scale & smaller scale & period 2 (alone) \\
\bottomrule
\end{tabular}
\caption{Direction of the equivalence-scale tilt in the Euler equation:
which period's total consumption gets the extra share, beyond what
$\beta R$ alone prescribes.}
\label{tab:tilt}
\end{table}

The estimated model has $\mu_{\text{low}} = \mu_{\text{high}} = 0.8 < 1$
(\texttt{est\_params\_alg1\_sparse.pkl}). Pre-fix, consumption always tilted
toward the \emph{larger}-household period, regardless of $\mu$. Post-fix
(current code), it tilts toward the \emph{smaller}-household period --
toward being alone -- because removing the multiplier is what made the
tilt direction $\mu$-dependent, and the estimated $\mu<1$ lands on the
reversed side.

\end{document}
